% ID: MAT-TRI-EQG-012
% Titolo: Intersezioni di una funzione goniometrica con gli assi (2)
% Disciplina: Matematica
% Area: Trigonometria
% Argomento: Equazioni goniometriche
% Classe: Quarta secondaria di secondo grado
% Difficolta: 4
% Tipo: grafico
% Tag: matematica, trigonometria, equazioni_goniometriche
% Fonte: verifica_goniometria_intermedia_anonimizzata
% Provenienza: materiale fornito dal docente; testo anonimizzato; soluzione redatta per Exergo

\begin{esercizio}
Trova le coordinate dei punti di intersezione con gli assi
cartesiani del grafico della funzione: 
\[
f(x)=\sin 2x+\sqrt{3}\cos 2x-1.
\]
Considera $x\in[0;2\pi)$.
Salvo diversa indicazione, lavora in $\mathbb R$; angoli in radianti.
\end{esercizio}

\begin{soluzione}
$2\sin(2x+\pi/3)=1$ dà $x=-\pi/12+k\pi$ o $x=\pi/4+k\pi$.
Nell'intervallo richiesto le ascisse sono $\pi/4,11\pi/12,5\pi/4,23\pi/12$:
le intersezioni sono $(a,0)$ per tali valori di $a$ e $(0,\sqrt3-1)$.
\end{soluzione}
